\scanpage{98}
évidente pour que $\Theta$ se relève en un $\widetilde{\Theta}$, c'est que l'image de l'application (51) soit \uline{stable} sous l'action de $\Theta$; en fait, l'action $\widetilde{\Theta}$ présomptive, restreinte à $E_0$, doit donner la compatibilité. \deleted{\unsure{Rien qu'en} \illegible{3} \unsure{Inversement}} \inserted{On voit donc qu'en deux pts de $\overline{\Sigma}$ conjugués sous $\Theta$, le degré de $\unsure{(52)|\overline{\Sigma}}$ (cardinal de la fibre) doit} \deleted{de stabilité} \deleted{\unsure{telle}} \unsure{qui par raison d'images signifie} justement qu'il existe une \deleted{$\unsure{\widetilde{\Theta}}$} \inserted{bijection}
\margin{\unsure{On lie} les deux conditions.}
\[
(52)\qquad \pi_0(\widetilde{\Theta}_0):\pi_0(E_0)\overset{\sim}{\longrightarrow}\pi_0(E_0)
\]
\deleted{telle que} qui soit compatible avec l'action de $\Theta$ sur le deuxième mb de (51). Ceci dit, pour \unsure{toute telle} relèvement, il existe bel \inserted{\unsure{p.e.}} et bien un $\widetilde{\Theta}$ qui l'induise, lequel $\widetilde{\Theta}$ est \inserted{$\widetilde{\Theta}_0:E_0\overset{\sim}{\to}E_0$ est compatible avec l'action de $\pi$ et l'action de $\Theta$ sur $\pi$ (cf (46))} \inserted{bien connu quand on} connaît ses valeurs $\widetilde{\Theta}(x_i)$ \deleted{\unsure{pour}} pour \uline{un} $x_i\in E_0$ dans chaque composante connexe de $E_0$ \ldots

Parmi les conditions nécessaires \inserted{importantes} pour un automorphisme $\Theta$ \unsure{d'un} $\widehat{\Pi}_1(X,I)$, $I\subset X_K$, pour provenir de $\Gamma$, j'ai oublié apparemment la plus importante de toutes, avec laquelle je travaille \deleted{\unsure{pourtant}} constamment --- à tel point que cela finit par sembler aller de soi!
\mspage{97}
